\section*{08. The disk identity as turning number}
\begin{center}\begin{minipage}{.86\textwidth}\centering
\begin{tikzpicture}[bf,x=1.05cm,y=1.05cm]
\path[use as bounding box] (-5.80,-5.80) rectangle (5.80,5.80);
\fill[bfblack,even odd rule] (0,0) circle[radius=5.55] (0,0) circle[radius=5.0];
\fill[bfblack] (-4.330,-2.500) -- (-2.133,-0.405) -- (-2.200,3.528) -- (-4.330,-2.500) -- cycle;
\fill[bfblack] (-4.330,-2.500) -- (-0.344,-4.394) -- (-1.614,-2.366) -- (-4.330,-2.500) -- cycle;
\fill[bfblack] (-2.133,-0.405) -- (0.764,-0.167) -- (-0.963,1.243) -- (-2.133,-0.405) -- cycle;
\fill[bfblack] (-2.133,-0.405) -- (-1.614,-2.366) -- (0.900,-2.302) -- (-2.133,-0.405) -- cycle;
\fill[bfblack] (-0.344,-4.394) -- (4.330,-2.500) -- (0.900,-2.302) -- (-0.344,-4.394) -- cycle;
\fill[bfblack] (4.330,-2.500) -- (2.799,2.992) -- (2.506,-0.651) -- (4.330,-2.500) -- cycle;
\fill[bfblack] (0.000,5.000) -- (-2.200,3.528) -- (2.799,2.992) -- (0.000,5.000) -- cycle;
\fill[bfblack] (-2.200,3.528) -- (-0.963,1.243) -- (0.610,2.065) -- (-2.200,3.528) -- cycle;
\fill[bfblack] (0.764,-0.167) -- (2.799,2.992) -- (0.610,2.065) -- (0.764,-0.167) -- cycle;
\fill[bfblack] (0.764,-0.167) -- (0.900,-2.302) -- (2.506,-0.651) -- (0.764,-0.167) -- cycle;
\fill[bfR!12!bfblack] (-2.133,-0.405) -- (0.764,-0.167) -- (-0.963,1.243) -- (-2.133,-0.405) -- cycle;
\fill[bfR!12!bfblack] (-2.133,-0.405) -- (-1.614,-2.366) -- (0.900,-2.302) -- (-2.133,-0.405) -- cycle;
\fill[bfR!12!bfblack] (-2.200,3.528) -- (-0.963,1.243) -- (0.610,2.065) -- (-2.200,3.528) -- cycle;
\fill[bfR!12!bfblack] (0.764,-0.167) -- (2.799,2.992) -- (0.610,2.065) -- (0.764,-0.167) -- cycle;
\fill[bfR!12!bfblack] (0.764,-0.167) -- (0.900,-2.302) -- (2.506,-0.651) -- (0.764,-0.167) -- cycle;
\fill[bfR!7] (-2.133,-0.405) -- (-0.963,1.243) -- (-2.200,3.528) -- (-2.133,-0.405) -- cycle;
\fill[bfR!7] (-2.133,-0.405) -- (0.900,-2.302) -- (0.764,-0.167) -- (-2.133,-0.405) -- cycle;
\fill[bfR!7] (-1.614,-2.366) -- (-0.344,-4.394) -- (0.900,-2.302) -- (-1.614,-2.366) -- cycle;
\fill[bfR!7] (-2.200,3.528) -- (0.610,2.065) -- (2.799,2.992) -- (-2.200,3.528) -- cycle;
\fill[bfR!7] (-0.963,1.243) -- (0.764,-0.167) -- (0.610,2.065) -- (-0.963,1.243) -- cycle;
\fill[bfR!7] (0.764,-0.167) -- (2.506,-0.651) -- (2.799,2.992) -- (0.764,-0.167) -- cycle;
\draw[bf edge] (0,0) circle[radius=5.0];
\draw[bf edge] (-4.330,-2.500) -- (-2.133,-0.405);
\draw[bf edge] (-4.330,-2.500) -- (-1.614,-2.366);
\draw[bf edge] (-4.330,-2.500) -- (-0.344,-4.394);
\draw[bf edge] (-4.330,-2.500) -- (-2.200,3.528);
\draw[bf edge] (-2.133,-0.405) -- (-1.614,-2.366);
\draw[bf edge] (-2.133,-0.405) -- (-2.200,3.528);
\draw[bf edge] (-2.133,-0.405) -- (-0.963,1.243);
\draw[bf edge] (-2.133,-0.405) -- (0.764,-0.167);
\draw[bf edge] (-2.133,-0.405) -- (0.900,-2.302);
\draw[bf edge] (-1.614,-2.366) -- (-0.344,-4.394);
\draw[bf edge] (-1.614,-2.366) -- (0.900,-2.302);
\draw[bf edge] (-0.344,-4.394) -- (4.330,-2.500);
\draw[bf edge] (-0.344,-4.394) -- (0.900,-2.302);
\draw[bf edge] (4.330,-2.500) -- (0.900,-2.302);
\draw[bf edge] (4.330,-2.500) -- (2.506,-0.651);
\draw[bf edge] (4.330,-2.500) -- (2.799,2.992);
\draw[bf edge] (0.000,5.000) -- (-2.200,3.528);
\draw[bf edge] (0.000,5.000) -- (2.799,2.992);
\draw[bf edge] (-2.200,3.528) -- (-0.963,1.243);
\draw[bf edge] (-2.200,3.528) -- (2.799,2.992);
\draw[bf edge] (-2.200,3.528) -- (0.610,2.065);
\draw[bf edge] (-0.963,1.243) -- (0.764,-0.167);
\draw[bf edge] (-0.963,1.243) -- (0.610,2.065);
\draw[bf edge] (0.764,-0.167) -- (0.900,-2.302);
\draw[bf edge] (0.764,-0.167) -- (2.506,-0.651);
\draw[bf edge] (0.764,-0.167) -- (2.799,2.992);
\draw[bf edge] (0.764,-0.167) -- (0.610,2.065);
\draw[bf edge] (0.900,-2.302) -- (2.506,-0.651);
\draw[bf edge] (2.506,-0.651) -- (2.799,2.992);
\draw[bf edge] (2.799,2.992) -- (0.610,2.065);
\draw[bf D] (-2.133,-0.405) -- (-1.614,-2.366);
\draw[bf D] (-1.614,-2.366) -- (-0.344,-4.394);
\draw[bf D] (-0.344,-4.394) -- (0.900,-2.302);
\draw[bf D] (0.900,-2.302) -- (2.506,-0.651);
\draw[bf D] (2.506,-0.651) -- (2.799,2.992);
\draw[bf D] (2.799,2.992) -- (-2.200,3.528);
\draw[bf D] (-2.200,3.528) -- (-2.133,-0.405);
\node[fill=white] at (-2.304,-0.437) {$-3$};
\node[fill=white] at (-1.743,-2.555) {$3$};
\node[fill=white] at (-0.372,-4.746) {$2$};
\node[fill=white] at (0.972,-2.486) {$-3$};
\node[fill=white] at (2.706,-0.703) {$3$};
\node[fill=white] at (3.023,3.231) {$2$};
\node[fill=white] at (-2.376,3.810) {$2$};
\draw[bf root] (-4.330,-2.500) circle[radius=2.3pt];
\fill[bf vertex] (-2.133,-0.405) circle[radius=1.4pt];
\fill[bf vertex] (-1.614,-2.366) circle[radius=1.4pt];
\fill[bf vertex] (-0.344,-4.394) circle[radius=1.4pt];
\draw[bf root] (4.330,-2.500) circle[radius=2.3pt];
\draw[bf root] (0.000,5.000) circle[radius=2.3pt];
\fill[bf vertex] (-2.200,3.528) circle[radius=1.4pt];
\fill[bf vertex] (-0.963,1.243) circle[radius=1.4pt];
\fill[bf vertex] (0.764,-0.167) circle[radius=1.4pt];
\fill[bf vertex] (0.900,-2.302) circle[radius=1.4pt];
\fill[bf vertex] (2.506,-0.651) circle[radius=1.4pt];
\fill[bf vertex] (2.799,2.992) circle[radius=1.4pt];
\fill[bf vertex] (0.610,2.065) circle[radius=1.4pt];
\end{tikzpicture}
\end{minipage}\end{center}
The labels are signed exterior turns in units of pi over three, computed from the two incident directed boundary edges. Their sum is six on this counterclockwise cycle. Give white corners angle one, the chosen black corner angle five, and the other black corners angle minus one, all in the same units. Each bounded face has angle sum three. Every nonroot vertex has angle sum six because it is chosen exactly once. Euler counting then gives the disk a total boundary turn of six. Negative corner angles are allowed; this is a combinatorial angle assignment, not a Euclidean embedding with ordinary positive angles. The root face is outside the disk, where its exceptional angle sum cannot interfere. The topology referee supplied the boundary flatness argument needed to handle a clockwise cycle as well.

\textbf{Status.} PROVED-HERE in the topology submission, CORRECT after explanatory additions by its independent AI referee. COMPUTED turn labels here; human review pending.

\textbf{Connects to.} 6 Euler census; 7 white surplus; 21 hypotheses.

{\footnotesize\raggedright Sources: research/\allowbreak{}2026-11-14\_opus\_route-topology/\allowbreak{}TOPOLOGY.md, Section 2;\allowbreak{} research/\allowbreak{}2026-11-14\_fable\_referee-topology/\allowbreak{}REFEREE.md, part a.\par}
